\documentclass{abntex2}
\usepackage[utf8]{inputenc}
\usepackage[T1]{fontenc}
\usepackage{geometry}
\usepackage{graphicx}
\usepackage{float}
\usepackage{hyperref}
\usepackage[english,indonesian]{babel}

\title{Optimasi Pemrosesan Template LaTeX Menggunakan Pendekatan Berbasis AI}
\author{Academic Support Agent}
\date{\today}

\begin{document}
\maketitle

\begin{abstract}

\end{abstract}

\chapter{PENDAHULUAN}


\chapter{TINJAUAN PUSTAKA}


% Original content (commented out):
% % Modelo de relatório de experimento em formato de relatório técnico

% Criado a partir de "Modelo de Relatório abnTeX", de autoria de Gustavo Yudi Bientinezi Matsuzake, disponibilizado na Galeria Overleaf (https://pt.overleaf.com/latex/templates/modelo-de-relatorio-abntex/cffmvqrhsprn).

% ------------------------------------------------------------------------

\documentclass[
	% -- opções da classe memoir --
	12pt,				% tamanho da fonte
	openright,			% capítulos começam em pág ímpar (insere página vazia caso preciso)
	twoside,			% para impressão em recto e verso. Oposto a oneside
	a4paper,			% tamanho do papel. 
	% -- opções da classe abntex2 --
	%chapter=TITLE,		% títulos de capítulos convertidos em letras maiúsculas
	%section=TITLE,		% títulos de seções convertidos em letras maiúsculas
	%subsection=TITLE,	% títulos de subseções convertidos em letras maiúsculas
	%subsubsection=TITLE,% títulos de subsubseções convertidos em letras maiúsculas
	% -- opções do pacote babel --
	english,			% idioma adicional para hifenização
	french,				% idioma adicional para hifenização
	spanish,			% idioma adicional para hifenização
	brazil,				% o último idioma é o principal do documento
	]{abntex2}


% ---
% PACOTES
% ---

\input{conf/pacotes.tex}

% --- 
% CONFIGURAÇÕES DE PACOTES
% --- 

\input{conf/preparativos.tex}

% ---
% Informações de dados para CAPA e FOLHA DE ROSTO
% ---
\titulo{COLOQUE SEU TÍTULO EM PORTUGUÊS AQUI, ELE PASSARÁ PARA A OUTRA LINHA AUTOMATICAMENTE SE PRECISO}
\autor{Nome do aluno 1 \\ Nome do aluno 2} % separe autores por \and ou, para saltar linha, por \\
\local{Brasília}
\data{2019}
\instituicao{%
  Universidade de Brasília -- UnB
  \par
  Instituto de Física
  \par
  Laboratório de Física 4 experimental}
\tipotrabalho{Relatório de experimento}
% O preambulo deve conter o tipo do trabalho, o objetivo, 
% o nome da instituição e a área de concentração 
\preambulo{Relatório de experimento apresentado como parte dos requisitos para aprovação na disciplina Física 4 experimental, turma A do período 2/2019, ministrado pelo professor Leonardo Luiz e Castro.}
% ---

% ---
% Configurações de aparência do PDF final


% alterando o aspecto da cor azul
\definecolor{blue}{RGB}{41,5,195}

% informações do PDF
\makeatletter
\hypersetup{
     	%pagebackref=true,
		pdftitle={\@title}, 
		pdfauthor={\@author},
    	pdfsubject={\imprimirpreambulo},
	    pdfcreator={LaTeX with abnTeX2},
		pdfkeywords={abnt}{latex}{abntex}{abntex2}{relatório técnico}, 
		colorlinks=true,       		% false: boxed links; true: colored links
    	linkcolor=blue,          	% color of internal links
    	citecolor=blue,        		% color of links to bibliography
    	filecolor=magenta,      		% color of file links
		urlcolor=blue,
		bookmarksdepth=4
}
\makeatother
% --- 

% --- 
% Espaçamentos entre linhas e parágrafos 
% --- 

% O tamanho do parágrafo é dado por:
\setlength{\parindent}{1.3cm}

% Controle do espaçamento entre um parágrafo e outro:
\setlength{\parskip}{0.2cm}  % tente também \onelineskip

% ---
% compila o indice
% ---
\makeindex
% ---

% ----
% Início do documento
% ----
\begin{document}

% Seleciona o idioma do documento (conforme pacotes do babel)
%\selectlanguage{english}
\selectlanguage{brazil}

% Retira espaço extra obsoleto entre as frases.
\frenchspacing 

% ----------------------------------------------------------
% ELEMENTOS PRÉ-TEXTUAIS
% ----------------------------------------------------------
% \pretextual

% ---
% Capa
% ---
\imprimircapa
% ---

% ---
% Folha de rosto
% (o * indica que haverá a ficha bibliográfica)
% ---
\imprimirfolhaderosto*
% ---

% ---
% Anverso da folha de rosto:
% ---

\input{textos/anverso.tex}

% ---
% Agradecimentos
% ---

\begin{agradecimentos}
    
    \input{textos/agradecimentos.tex}

\end{agradecimentos}

% ---

% ---
% RESUMO
% ---

% resumo na língua vernácula (obrigatório)
\setlength{\absparsep}{18pt} % ajusta o espaçamento dos parágrafos do resumo
\begin{resumo}

    \input{textos/resumo.tex}

\end{resumo}
% ---

% ---
% inserir lista de ilustrações
% ---
\pdfbookmark[0]{\listfigurename}{lof}
\listoffigures*
\cleardoublepage
% ---

% ---
% inserir lista de tabelas
% ---
\pdfbookmark[0]{\listtablename}{lot}
\listoftables*
\cleardoublepage
% ---

% ---
% inserir lista de abreviaturas e siglas
% ---
\begin{siglas}

    \input{textos/siglas.tex}
  
\end{siglas}
% ---

% ---
% inserir lista de símbolos
% ---
\begin{simbolos}

    \input{textos/simbolos.tex}

\end{simbolos}
% ---

% ---
% inserir o sumario
% ---
\pdfbookmark[0]{\contentsname}{toc}
\tableofcontents*
\cleardoublepage
% ---


% ---
% ELEMENTOS TEXTUAIS
% ---
\textual

% ---
% Introdução
% ---

\input{textos/introducao.tex}

% Se quiser agregar vários capítulos em "partes", use o comando \part antes deles:

\input{textos/instrucoes.tex} % Apague esta linha no seu documento.

\part{Metodologia}

\input{textos/objetivos.tex}

\input{textos/materiais.tex}

\input{textos/procedimento.tex}

\part{Resultados}

\input{textos/dados.tex}

\input{textos/analise.tex}

\input{textos/discussao.tex}

% ---
% Finaliza a parte no bookmark do PDF
% para que se inicie o bookmark na raiz
% e adiciona espaço de parte no Sumário
% ---
\phantompart

% ---
% Conclusão
% ---
\chapter{Conclusão}
% ---

\input{textos/conclusao.tex}

% ---
% ELEMENTOS PÓS-TEXTUAIS
% ---
\postextual

% ---
% Referências bibliográficas
% ---

\bibliography{relatorio}

% ---
% Glossário
% ---
%
% Consulte o manual da classe abntex2 para orientações sobre o glossário.
%
%\glossary

% ---
% Apêndices
% ---

% ---
% Inicia os apêndices
% ---
\begin{apendicesenv}

% Imprime uma página indicando o início dos apêndices
\partapendices

\input{textos/apendice-a.tex}

\input{textos/apendice-b.tex}

\end{apendicesenv}
% ---


% ----------------------------------------------------------
% Anexos
% ----------------------------------------------------------

% ---
% Inicia os anexos
% ---
\begin{anexosenv}

% Imprime uma página indicando o início dos anexos
\partanexos

\input{textos/anexo-a.tex}

\input{textos/anexo-b.tex}

\end{anexosenv}

%---------------------------------------------------------------------
% INDICE REMISSIVO
%---------------------------------------------------------------------

\phantompart

\printindex

%---------------------------------------------------------------------
% Formulário de Identificação (opcional)
%---------------------------------------------------------------------

\input{textos/formulario-identificacao.tex}

\end{document}


\end{document}